\begin{titlepage}
    \centering
    \thispagestyle{empty}

    \vspace*{2.0cm}

    {\LARGE\bfseries Trabajo Final de Overleaf\par}

    \vspace{0.6cm}
    \rule{0.75\textwidth}{0.6pt}
    \vspace{1.2cm}

    {\Large Estabilidad térmica versus eficiencia catalítica\par}

    \vspace{0.45cm}

    {\large\itshape
    ¿Puede la ingeniería enzimática mejorar ambas propiedades
    en biocatalizadores industriales?\par}

    \vspace{1.2cm}
    \rule{0.55\textwidth}{0.4pt}

    \vfill

    {\large Andy Edu Villalobos Melendez\par}
    \vspace{0.35cm}
    {\large Pontificia Universidad Católica del Perú\par}
    \vspace{1.2cm}
    {\large \today\par}

    \vspace*{1.2cm}
\end{titlepage}

\tableofcontents

\vspace{0.5cm}

\noindent
\hyperref[sec:introduccion]{Ir a Introducción}
\hfill
\hyperref[sec:metodologia]{Ir a Metodología}
\hfill
\hyperref[sec:evidencia]{Ir a Resultados}

\newpage

\section{Introducción y Justificación}
\label{sec:introduccion}

\subsection{El Problema Industrial}
\label{subsec:problema_industrial}

Las enzimas son proteínas que aceleran reacciones químicas y pueden utilizarse
como biocatalizadores para transformar materias primas en productos de interés.
Su aplicación incluye alimentos, detergentes, biocombustibles, productos
farmacéuticos y compuestos de alto valor.

Para una empresa no basta con que una enzima sea rápida. Debe conservar su
capacidad de conversión bajo las condiciones reales de operación. Temperaturas
elevadas, altas concentraciones de sustrato, acumulación de producto o periodos
prolongados de uso pueden reducir su desempeño
\citep{siddiqui_defying_2017}.

Por ello, la decisión industrial parte de dos preguntas:

\begin{enumerate}
    \item ¿Con qué eficiencia transforma la materia prima?
    \item ¿Durante cuánto tiempo conserva esa capacidad?
\end{enumerate}

Una enzima muy activa pero de corta duración puede perder rápidamente su ventaja
inicial. Una enzima muy resistente pero poco eficiente puede durar más sin
generar suficiente producto.

\subsection{El Conflicto Actividad--Estabilidad}
\label{subsec:tradeoff}

La estabilidad térmica favorece que la enzima conserve su estructura, mientras
que la catálisis requiere cierta movilidad molecular para reconocer el sustrato,
acomodarlo en el sitio activo y liberar el producto.

Una modificación que rigidice demasiado la proteína puede aumentar estabilidad
y, al mismo tiempo, limitar movimientos necesarios para la reacción. Esta
tensión se conoce como trade-off actividad--estabilidad
\citep{hou_tradeoff_2023}.

Desde una perspectiva industrial existen cuatro escenarios:

\begin{enumerate}
    \item alta eficiencia y baja estabilidad: produce rápido, pero durante poco tiempo;
    \item baja eficiencia y alta estabilidad: dura más, pero convierte lentamente;
    \item baja eficiencia y baja estabilidad: ofrece poco valor para el proceso;
    \item alta eficiencia y alta estabilidad: combina velocidad y duración.
\end{enumerate}

El cuarto escenario es alcanzable: se han reportado variantes enzimáticas en las
que estabilidad y actividad mejoran simultáneamente
\citep{siddiqui_defying_2017,hou_tradeoff_2023}.

\subsection{Por Qué Importa Analizar Ambas Propiedades}
\label{subsec:justificacion}

El análisis resulta relevante por tres razones:

\begin{enumerate}
    \item Razón productiva. Una mayor vida operativa puede sostener la conversión
    y elevar el producto obtenido en una misma campaña.

    \item Razón tecnológica. Evolución dirigida, diseño racional y herramientas
    computacionales permiten modificar regiones específicas de una proteína
    \citep{siddiqui_defying_2017}.

    \item Razón de decisión. Mejorar una propiedad no garantiza mejorar el
    proceso: también intervienen temperatura, sustrato, producto acumulado y
    selectividad.
\end{enumerate}

La propiedad molecular, por tanto, debe evaluarse por su efecto sobre el
resultado productivo y no únicamente por su valor aislado.

\subsection{Pregunta y Criterio de Análisis}
\label{subsec:pregunta}

El trabajo responde la siguiente pregunta:

\begin{center}
    \textit{¿El aumento de la termoestabilidad de una enzima industrial
    necesariamente reduce su eficiencia catalítica?}
\end{center}

Para responderla se consideran tres niveles:

\begin{enumerate}
    \item termoestabilidad: cuánto tiempo conserva actividad;
    \item eficiencia catalítica: qué tan eficazmente convierte el sustrato;
    \item desempeño productivo: cuánto producto genera durante la operación.
\end{enumerate}

La lógica central es:

\begin{equation}
\text{Estabilidad}
+
\text{Eficiencia}
+
\text{Condiciones de Operación}
\longrightarrow
\text{Producción Útil}.
\end{equation}

Las variables se definen en la \hyperref[sec:variables]{Sección 2}, el protocolo
en la \hyperref[sec:metodologia]{Sección 3}, los resultados en la
\hyperref[sec:evidencia]{Sección 4} y la decisión industrial en la
\hyperref[sec:discusion]{Sección 5}.
